\documentclass[10pt,a4paper]{amsart}
\usepackage[margin=19mm]{geometry}
\usepackage{amsmath,amssymb,mathtools}
\usepackage[scr=rsfs,cal=euler]{mathalfa}
\usepackage{xfrac}
\usepackage[unicode,hidelinks,bookmarks=false]{hyperref}
\newcommand{\ie}{i.e.\ }
\newcommand{\cf}{cf.\ }
\newcommand{\resp}{respectively\ }
\newcommand{\Subsection}{\subsection}
\providecommand{\nobreakdash}{\nobreak\hbox{-}}
\setlength{\emergencystretch}{2em}
\multlinegap0pt
\usepackage{amssymb}
\usepackage{mathtools}
\usepackage{xcolor}
\newtheorem{lemma}{Lemma}
\title{Affine Copies of Three-Point Patterns in Sets of Integers}
\author{Samuel Korsky}
\date{Source: arXiv:2609.02308v1 (CC BY 4.0)}
\begin{document}
\maketitle

\noindent\emph{Source and licence.} Affine Copies of Three-Point Patterns in Sets of Integers. Version arXiv:2609.02308v1 (CC BY 4.0), distributed by arXiv under the Creative Commons Attribution 4.0 International licence, http://creativecommons.org/licenses/by/4.0/, as recorded in the arXiv licence metadata for this exact version and shown on the abstract page https://arxiv.org/abs/2609.02308v1. Full licence notice: \texttt{licenses/arxiv-2609.02308-CC-BY-4.0.txt}.\par\smallskip

\noindent\textit{Editorial scope.} Counted unit: the article's lemma lem:one-three-way-split (source0.tex lines 1297--1307). The article's complete original proof of it is delivered with it (source0.tex lines 1308--1473, one \texttt{proof} environment of 166 lines). No mathematics is written, repaired or re-derived: the statement and the proof are the article's own lines, in the article's own notation, and no retained line is substituted or re-flowed.  No further source-local context is needed: every cross-reference of every retained line resolves inside the two delivered spans.  Nothing was omitted from any retained span, so the package carries no omission record.  No retained line contains a citation, so the wrapper carries no bibliography and no bibliography entry.  No retained line contains a figure environment, an image-inclusion command or drawing code, so no image is delivered and the package declares no asset dependency.  Editorial formatting only, in addition: an AMS \texttt{amsart} preamble that restates the article's own declarations and macros used by the retained lines, namely the environments \texttt{lemma} on separate counters, together with the standard packages those lines need.  The retained environments print in this document as Lemma 1 in order of appearance, whereas the article prints the same items with the numbering of its own full document.  The complete native source of the article as distributed in the arXiv e-print for this exact version is bundled unchanged as \texttt{sources/209190/source0.tex}.
\begin{lemma}[The inequality needed at one split]
\label{lem:one-three-way-split}
For every $\mathbf p=(p_0,p_1,p_2)$ with $p_i\ge0$ and $\sum_i p_i=1$,
\begin{equation}\label{eq:one-three-way-split}
 V(\mathbf p)+\sum_{i=0}^2p_i^2
 K\left(\frac{p_j}{p_i},\frac{p_k}{p_i}\right)
 \le c\left(1-\sum_{i=0}^2p_i^2\right),
\end{equation}
where $\{i,j,k\}=\{0,1,2\}$.  If $p_i=0$, the corresponding term is
understood as its limiting value.
\end{lemma}

\begin{proof}
The proof is an explicit quadratic calculation.  We first sort the three
parts and express their ratios using two variables.  The maximum with zero in $V$ then gives two cases.  One case follows from a simple uniform bound.  In
the other, the relevant functions are quadratic on three intervals; after
clearing denominators, only four expressions remain to be checked.
Set
\begin{equation}\label{eq:def-H-g}
 H(t)=F(t)-\eta\beta(t),
 \qquad
 g(t)=\beta(t)-1.
\end{equation}
Then
\begin{equation}\label{eq:K-Hg}
 K(u,v)=H(u)+H(v)+\eta(g(u)+g(v))_+.
\end{equation}
For $t\ge1$, we have $F(t)=1/3$, $\beta(t)=1$, and hence
$H(t)=\rho$ and $g(t)=0$.  For every $t\ge0$,
\begin{equation}\label{eq:single-rho}
 H(t)+\eta(g(t))_+
 =F(t)-\eta\min\{\beta(t),1\}\le\rho.
\end{equation}
Indeed, if $\beta(t)\ge1$, use $F(t)\le1/3$.  If $\beta(t)<1$, then
$t<1/7$, and the left side is at most
$F(1/7)=25/196<\rho$.

By symmetry, assume $p_0\ge p_1\ge p_2$.  The cases with a zero coordinate
are obtained by taking $r\downarrow0$ when $p_2=0<p_1$, or
$x\downarrow0$ when $p_1=p_2=0$, in the parameterization below.  Terms with
the corresponding squared prefactors vanish, and the remaining terms have
finite one-sided limits.  It therefore suffices to write
\begin{equation}\label{eq:param-p}
 \mathbf p=\frac{(1,x,xr)}d,
 \qquad d=1+x+xr,
 \qquad 0<x,r\le1,
\end{equation}
and put
\begin{equation}\label{eq:def-Q}
 Q=x+xr+x^2r.
\end{equation}
Thus $2Q/d^2=1-\sum_i p_i^2$.  After multiplying
\eqref{eq:one-three-way-split} by $d^2$, the right side minus the left side is
\begin{equation}\label{eq:def-Gamma}
 \Gamma=2cQ-2(Q-\rho d^2)_+
 -K(x,xr)-x^2K(1/x,r)-2\rho x^2r^2.
\end{equation}
We prove $\Gamma\ge0$.

First suppose $Q\ge\rho d^2$.  Equations \eqref{eq:K-Hg} and
\eqref{eq:single-rho}, together with
$(a+b)_+\le a_++b_+$, give
\[
 K(x,xr)\le2\rho,
 \qquad
 K(1/x,r)\le2\rho.
\]
In this case \eqref{eq:def-Gamma} simplifies to
\[
 \Gamma=(2\rho-K(x,xr))
 +x^2(2\rho-K(1/x,r))\ge0.
\]
Now suppose $Q\le\rho d^2$.  Define two auxiliary functions
\begin{equation}\label{eq:def-RS}
 R(t)=732\bigl(2ct-\rho t^2-H(t)\bigr),
 \qquad
 S(t)=R(t)-19g(t).
\end{equation}
The factor $732$ clears the denominators, and the two functions record the
two possibilities created by the positive part in \eqref{eq:K-Hg}: $R$ is
the inactive case, while $S=R-19g$ is the active case.  Since
$732\eta=19$, direct substitution gives
\begin{equation}\label{eq:min-RS}
 732\Gamma
 =\min\{R(x)+R(xr),S(x)+S(xr)\}
 +x^2\min\{R(r),S(r)\}.
\end{equation}
Introduce the three quadratic polynomials
\begin{equation}\label{eq:def-ABC}
 A(t)=3t+58t^2,
 \qquad
 B(t)=(6t-1)(54t-19),
 \qquad
 C(t)=-3(75t^2-188t+75).
\end{equation}
Substituting the definitions of $F$ and $\beta$ yields
\begin{equation}\label{eq:R-pieces}
 R(t)=
 \begin{cases}
 A(t),&0\le t\le3/8,\\
 B(t),&3/8<t\le2/3,\\
 C(t),&2/3<t\le1,
 \end{cases}
 \qquad
 S(t)=
 \begin{cases}
 B(t),&0\le t\le2/3,\\
 C(t),&2/3<t\le1.
 \end{cases}
\end{equation}
These formulas imply
\begin{equation}\label{eq:RS-basic}
 R(t)\ge0,
 \qquad
 R(t)\ge\frac{25}{9}t^2,
 \qquad
 S(t)\ge-\frac{25}{9}
 \qquad(0\le t\le1).
\end{equation}
For the last bound, note that
\begin{equation}\label{eq:B-square}
 B(t)=324\left(t-\frac7{27}\right)^2-\frac{25}{9}.
\end{equation}
Also $S(t)=C(t)\ge51$ when $t>2/3$.

Expanding the two minima in \eqref{eq:min-RS} gives four possible sums.
Two follow immediately from \eqref{eq:RS-basic}:
\[
 R(x)+R(xr)+x^2R(r)\ge0
\]
and
\[
 R(x)+R(xr)+x^2S(r)
 \ge R(x)-\frac{25}{9}x^2\ge0.
\]
For the third sum, we prove
\begin{equation}\label{eq:SS-goal}
 S(x)+S(xr)+x^2S(r)\ge0.
\end{equation}
If $x>2/3$, the bound $S(x)\ge51$ and
$S\ge-25/9$ make this immediate.  Suppose $x\le2/3$, so also
$xr\le2/3$.  If $r\le2/3$, all three terms use the polynomial $B$, and
\begin{equation}\label{eq:critical-sos}
 \begin{aligned}
 B(x)+B(xr)+x^2B(r)
 ={}&648\left(xr-\frac{7(x+1)}{54}\right)^2\\
 &+\frac{61}{9}(7x-2)^2\ge0.
 \end{aligned}
\end{equation}
If $r>2/3$, put $a(r)=99r^2+564r+99$.  Then
\begin{equation}\label{eq:SS-large-r}
 \begin{aligned}
 B(x)+B(xr)+x^2C(r)
 ={}&a(r)\left(x-\frac{84(1+r)}{a(r)}\right)^2\\
 &+\frac{366(20r-9r^2-9)}{a(r)}\ge0.
 \end{aligned}
\end{equation}
The final numerator is positive on $[2/3,1]$.

The remaining sum is
\begin{equation}\label{eq:SR-goal}
 S(x)+S(xr)+x^2R(r).
\end{equation}
If $r\ge1/7$, then $g(r)\ge0$, so $R(r)\ge S(r)$, and
\eqref{eq:SS-goal} proves the claim.  Suppose $r<1/7$.  If $x>2/3$,
then $S(x)\ge51$, while $S(xr)>0$ and $R(r)\ge0$.  If $x\le2/3$, put
$d(r)=382r^2+3r+324$.  Here $S(x)=S(xr)=B$ and $R(r)=A$, and
\begin{equation}\label{eq:SR-small-r}
 \begin{aligned}
 B(x)+B(xr)+x^2A(r)
 ={}&d(r)\left(x-\frac{84(1+r)}{d(r)}\right)^2\\
 &+\frac{2(3730r^2-6999r+2628)}{d(r)}\ge0.
 \end{aligned}
\end{equation}
The final quadratic decreases and remains positive on $[0,1/7]$.
Thus all four possible sums in \eqref{eq:min-RS} are nonnegative, proving
\eqref{eq:one-three-way-split}.
\end{proof}

\end{document}
